\documentclass[journal]{IEEEtran}

% --- Paquetes ---
\usepackage[spanish]{babel}
\usepackage{csquotes}
\usepackage[backend=biber,style=ieee]{biblatex}
\addbibresource{ref.bib}
\usepackage{lipsum} 

\usepackage{fontspec}
\setmainfont{TeX Gyre Heros} % Fuente sans-serif para todo el texto
\setsansfont{TeX Gyre Heros}
\setmonofont{TeX Gyre Cursor} % Para el código fuente

% --- Título y autores (EN EL PREÁMBULO) ---
\title{TÍTULO DEL PROYECTO}

\author{
  \IEEEauthorblockN{Nombre 1 \and Nombre 2 \\[1.5ex]
  {\normalsize Universidad del Quindío}\\
  {\normalsize Facultad de Ingeniería}\\
  {\normalsize Ingeniería de Sistemas y Computación}\\
  {\normalsize NOMBRE DE LA ASIGNATURA}}
}

\begin{document}

\maketitle

\begin{abstract}
Resumen del documento (150--250 palabras). \lipsum[1]
\end{abstract}

\begin{IEEEkeywords}
palabra clave, palabra clave, palabra clave
\end{IEEEkeywords}

\section{Introducción}
\lipsum[2]

No olvidar usar referencias \parencite{name}.

\section{NOMBRE DE LA PRIMERA SECCIÓN}
\lipsum[3]

\subsection{Nombre de la sub-sección}
\lipsum[4]

\section{NOMBRE DE LA SEGUNDA SECCIÓN}
\lipsum[5]

\subsection{Nombre de la sub-sección}
\lipsum[6]

\section{Conclusiones}
\lipsum[8]

\printbibliography

\end{document}
